\scanpage{57}

\begin{center}- 155 -\end{center}
\begin{flushright}IX\end{flushright}

comme le conoyau de l'homomorphisme $\mathbb Z^{I\times J}\longrightarrow\mathbb Z^K$ induit par $h$, lequel est transposé de $\mathbb Z^{(h)}$. Notons d'ailleurs que le graphe $\Gamma(C)$ s'explicite aussi de façon évidente en termes de $h$, par $\Gamma_{\unsure{o|0}}=J$, $\Gamma_1=(K\times_I K-\Delta(K))/\text{symétrie}$. D'autre part, choisissant une orientation sur chacune des arrêtes de $\Gamma(C)$, et désignant par $\Gamma_1$ l'ensemble de ces arrêtes $(x_i,C',C'')$, on a
\[
(12.3.11)\qquad
H_1(\Gamma(C),\mathbb Z)=\operatorname{Ker}(d:\mathbb Z^{(\Gamma_1)}\longrightarrow\mathbb Z^{(J)})
\]
où l'opérateur différentiel est donné par
\[
(12.3.12)\qquad d(x,C',C'')=g(C'')-g(C')\quad.
\]
Ceci posé, l'homomorphisme (12.3.8) est induit par l'homomorphisme $v$ de complexes de chaînes :
\[
(12.3.13)\quad
v=(v_{\unsure{o|0}},v_1):C_{\bullet}(\Gamma(C))
=[\mathbb Z^{(\Gamma)}\longrightarrow\mathbb Z^{(J)}]
\longrightarrow[\mathbb Z^{(K)}\longrightarrow\mathbb Z^{(I\times J)}]
\]
donné par
\[
(12.3.14)\quad
v_1(x,C',C'')=(x,C'')-(x,C'),\quad
v_{\unsure{o|0}}(\widetilde C_j)=(\unsure{o|0},\widetilde C_j)\quad.
\]

12.4. Nous nous plaçons maintenant dans le cas envisagé à la fin de 12.2. En vertu de (12.1.12) et de (12.3.6) appliqué à $C=X_{\unsure{o|0}}$, le groupe des caractères $M$ du tore maximal $T_{\unsure{o|0}}$ de la fibre spéciale $A_{\unsure{o|0}}$ du modèle de Néron $A$ de $A=\uline{\operatorname{Pic}}^{\unsure{o|0}}_{X_\eta/\eta}$ est donné par
\[
(12.4.1)\qquad M=\operatorname{Ker}(\mathbb Z^I\xrightarrow{d}\mathbb Z^J)\quad,
\]
où $I$ est l'ensemble des points singuliers de $X_{\unsure{o|0}}$, et $J$ l'ensemble de ses composantes irréductibles, l'homomorphisme $\mathbb Z^I\longrightarrow\mathbb Z^J$ étant donné par

\margin{473}
\margin{473}
\margin{56}
